\section{Related Work}
\label{sec:related}

This section reviews peer-reviewed studies published between 2022 and 2026 that are closest to
the present work. They are grouped into unimodal deep-learning approaches for the ECG, multimodal
approaches, federated-learning approaches and explainable approaches. The reported figures are
those stated by the respective authors on their own datasets and splits; they are therefore not
directly comparable with each other, and the metric is named whenever it is not accuracy.
Table~\ref{tab:related} summarises the review.

\subsection{Unimodal ML/DL Approaches}
Mehari and Strodthoff (2022) \cite{mehari2022ssl} proposed adapting contrastive predictive coding,
SimCLR, BYOL and SwAV to 12-lead ECGs from PTB-XL and achieved a macro AUROC of 94.18\% on the
71-statement task after self-supervised pretraining and fine-tuning, but the gain over supervised
training was about one percentage point and the pretraining relied on pooling several large public
corpora at one site. Na \textit{et al.} (2024) \cite{na2024stmem} proposed ST-MEM, a
spatio-temporal masked auto-encoder, using PTB-XL for evaluation on the five diagnostic
superclasses and achieved 82.5\% accuracy and a macro AUROC of 93.3\% after fine-tuning, but the
method requires large-scale centralised pretraining and does not explain its decisions. Mehari and
Strodthoff (2023) \cite{mehari2023s4} proposed structured state space (S4) models for the raw
12-lead signal and achieved a macro AUROC of 94.17\% on the PTB-XL 71-statement task from the
signal alone, but the model was trained on a single pooled dataset and its comparison with
convolutional baselines was limited to centralised training. Together these studies establish
strong single-site baselines, yet each assumes that all training records can be gathered in one
place.

\subsection{Multimodal Approaches}
Mehari and Strodthoff (2023) \cite{mehari2023s4} also proposed fusing basic demographic metadata
with the ECG representation of a pretrained S4 model on PTB-XL and achieved a macro AUROC of
94.63\%, but the metadata improved the result by only 0.18 percentage points and the fusion was
again centralised. Li \textit{et al.} (2023) \cite{li2023mets} proposed METS, which aligns ECG
signals with automatically generated reports through a frozen language model, using PTB-XL and
achieved 84.2\% zero-shot accuracy on the superclass task, but its macro F1 of 65.7\% remained below
the 77.6\% of a supervised ResNet-18 reported in the same study. Liu \textit{et al.} (2024)
\cite{liu2024merl} proposed MERL, an ECG--report multimodal pretraining method with test-time
clinical-knowledge prompts, using MIMIC-ECG for pretraining and six public datasets for evaluation
and achieved an average zero-shot AUROC of 75.2\%, but zero-shot performance stays well below
supervised fine-tuning and the method needs paired free-text reports for every training ECG.
Palakayala \textit{et al.} (2025) \cite{palakayala2025hamf} proposed HAMF, a hierarchical
attention-based fusion of MRI, gait and speech for Parkinson's disease, using PPMI, PhysioNet and
voice repositories and achieved 94.2\% accuracy, but the modalities were compiled from different
cohorts rather than acquired from the same patients and the model was trained centrally. These
studies show that additional modalities can help, but also that every modality must be available
for each patient at one site.

\subsection{Federated Learning Approaches}
Raza \textit{et al.} (2022) \cite{raza2022fed} proposed a federated transfer-learning framework with
a convolutional auto-encoder and an explanation module, using the MIT-BIH arrhythmia database and
achieved 98.9\% accuracy on clean signals (94.5\% on noisy signals), but the evaluation used
five-fold cross-validation on beats of a 48-record database rather than patient-level 12-lead
diagnosis. Goto \textit{et al.} (2022) \cite{goto2022fed} proposed multinational FL for ECG and
echocardiogram models across three hospitals in the United States and Japan and achieved AUROCs of
90\%--96\% on internal test sets and 93\% on an external site for hypertrophic cardiomyopathy, but
all sites trained the same architecture and the task was a single binary diagnosis. Quan
\textit{et al.} (2025) \cite{quan2025fedocw} proposed FedOcw, which optimises client-specific
aggregation weights in every round, using five multilingual speech corpora for Parkinson's disease
and achieved 74.81\% accuracy in the Spanish--Italian scenario, outperforming FedAvg, FedProx,
SCAFFOLD and FedNova, but the method still averages parameters of an identical network over 100
rounds. Sathi \textit{et al.} (2025) \cite{sathi2025xfl} proposed explainable FL on ECG fiducial
features with SHAP and LIME, using the MIT-BIH, European ST-T and Fantasia databases and achieved
92.3\% accuracy, but the clients shared one model family and the reported accuracies vary between
sections of the study. In all of these studies the clients are required to train the same model and
to communicate repeatedly.

\subsection{Explainable AI Approaches}
Jahmunah \textit{et al.} (2022) \cite{jahmunah2022gradcam} proposed DenseNet and CNN models with
Grad-CAM, using the PTB diagnostic database to separate healthy subjects from ten MI locations and
achieved accuracies above 95\%, but the data came from a small cohort and the heat maps were assessed
only qualitatively. Wagner \textit{et al.} (2024) \cite{wagner2024xai} proposed a set of building
blocks for auditing ECG models, comparing local attributions with concept-based global analyses on
PTB-XL, and reported that saliency was the only attribution method that passed all of their sanity
checks, but the study focuses on the explanation methods and does not propose a new classifier.
Pantelidis \textit{et al.} (2025) \cite{pantelidis2025xplaim} proposed ECG-XPLAIM, an inception-style
1D CNN with Grad-CAM trained on MIMIC-IV and externally validated on PTB-XL, and achieved a macro
AUROC of 98\% for conduction disturbances, but the model addresses selected arrhythmias and its
sensitivity for long-QT detection was only 69.1\%. Explainability is increasingly expected, yet it
is rarely combined with a privacy-preserving training scheme or with a mechanism that decides which
predictions should be referred to a clinician.

\begin{table*}[!t]
\centering
\caption{Summary of related work (reported figures are the authors' own; metric named when not accuracy)}
\label{tab:related}
\footnotesize
\setlength{\tabcolsep}{4pt}
\begin{tabular}{@{}p{0.6cm}p{2.4cm}p{3.1cm}p{3.6cm}p{2.0cm}p{4.8cm}@{}}
\toprule
Ref & Author \& Year & Dataset / Modality & Model / Method & Accuracy (\%) & Limitation \\
\midrule
\cite{mehari2022ssl} & Mehari \& Strodthoff, 2022 & PTB-XL; 12-lead ECG & CPC self-supervised pretraining + LSTM & 94.18 (macro AUROC) & Centralised pretraining on pooled corpora \\
\cite{na2024stmem} & Na \textit{et al.}, 2024 & PTB-XL superclasses; ECG & ST-MEM masked auto-encoder & 82.5 (AUROC 93.3) & Large centralised pretraining; no explanations \\
\cite{mehari2023s4} & Mehari \& Strodthoff, 2023 & PTB-XL; ECG (+ metadata) & S4 state space model & 94.17 / 94.63 (macro AUROC) & Single site; metadata adds 0.18 points \\
\cite{li2023mets} & Li \textit{et al.}, 2023 & PTB-XL; ECG + report text & METS, frozen language model & 84.2 (zero-shot) & F1 65.7 below supervised ResNet-18 \\
\cite{liu2024merl} & Liu \textit{et al.}, 2024 & MIMIC-ECG + 6 datasets; ECG + text & MERL multimodal pretraining & 75.2 (avg.\ zero-shot AUROC) & Needs paired reports; below fine-tuning \\
\cite{palakayala2025hamf} & Palakayala \textit{et al.}, 2025 & PPMI, PhysioNet, voice; MRI, gait, speech & HAMF hierarchical attention fusion & 94.2 & Modalities from different cohorts; centralised \\
\cite{raza2022fed} & Raza \textit{et al.}, 2022 & MIT-BIH; single-lead beats & Federated transfer learning + XAI & 98.9 (clean), 94.5 (noisy) & Beat-level task; same model on all clients \\
\cite{goto2022fed} & Goto \textit{et al.}, 2022 & 4 hospitals; ECG, echo & Multinational FL & 90--96 (AUROC) & Binary task; homogeneous architectures \\
\cite{quan2025fedocw} & Quan \textit{et al.}, 2025 & 5 speech corpora & FedOcw optimised client weights & 74.81 & 100 rounds of parameter averaging \\
\cite{sathi2025xfl} & Sathi \textit{et al.}, 2025 & MIT-BIH, ST-T, Fantasia; fiducial features & Explainable FL with SHAP/LIME & 92.3 & Hand-crafted features; inconsistent reporting \\
\cite{jahmunah2022gradcam} & Jahmunah \textit{et al.}, 2022 & PTB; 12-lead ECG & DenseNet/CNN + Grad-CAM & $>$95 & Small cohort; qualitative heat maps \\
\cite{wagner2024xai} & Wagner \textit{et al.}, 2024 & PTB-XL; 12-lead ECG & Attribution sanity checks, concept analysis & -- (not reported) & Explanation study, no new classifier \\
\cite{pantelidis2025xplaim} & Pantelidis \textit{et al.}, 2025 & MIMIC-IV $\rightarrow$ PTB-XL; ECG & ECG-XPLAIM 1D CNN + Grad-CAM & 98 (macro AUROC, CD) & Selected arrhythmias; LQT sensitivity 69.1 \\
\midrule
This work & -- & PTB-XL superclasses; 12-lead ECG & One-shot federated gate over 4 heterogeneous clients & 86.9 (label-wise); AUROC 90.5 & See Section~\ref{sec:discussion} \\
\bottomrule
\end{tabular}
\end{table*}

\subsection{Research Gap}
The reviewed studies reveal three gaps that motivate this work. Single-site ECG models achieve high
AUROC but assume that all records can be pooled. Federated methods remove this assumption but
average parameters, so every institution must adopt the same architecture and participate in many
synchronised rounds. Explainable models, finally, are seldom combined with privacy-preserving
training or with an explicit rule for deferring uncertain findings. The proposed model addresses
these gaps by letting each institution train its own architecture once, combining only calibrated
probabilities through a validated simplex gate, explaining the combined prediction with
attributions on the raw signal, and reporting automatically only those findings for which the
validated precision and recall reach the clinical target.
